\section{Taylor series, logarithmic derivatives, and power-series certificates}\label{sec:taylor}

The canonical-factor calculation \cref{eq:qseries} contains the full Taylor
series.  On \(|z|<1\), absolute convergence permits the order of summation to be
reversed:
\begin{align}
 \sum_{k=1}^{\infty}q_k(z)
 &=\sum_{k=1}^{\infty}\sum_{m=3}^{\infty}
   \frac{(-1)^{m+1}z^m}{m k^{m-1}}\notag\\
 &=\sum_{m=3}^{\infty}
   \frac{(-1)^{m+1}\zeta(m-1)}{m}z^m.
 \label{eq:q-double-sum}
\end{align}
The constructive justification is a rectangular-tail estimate.  If
\(|z|\le\rho<1\), then \(\zeta(m-1)<2\) for every \(m\ge3\), and hence all
omitted terms are dominated by a geometric series.

\begin{theorem}[Taylor series for \(\log G\)]\label{thm:taylor}
For \(|z|<1\), on the branch normalized by \(\log G(1)=0\),
\begin{equation}\label{eq:taylor-logG}
 \log G(1+z)
 =\frac{\log(2\pi)-1}{2}z
  -\frac{1+\gamma}{2}z^2
  +\sum_{m=3}^{\infty}
    \frac{(-1)^{m-1}\zeta(m-1)}{m}z^m.
\end{equation}
Equivalently, with \(k=m-1\),
\begin{equation}\label{eq:taylor-wiki}
 \log G(1+z)
 =\frac z2\log(2\pi)
  -\frac{z+(1+\gamma)z^2}{2}
  +\sum_{k=2}^{\infty}
    (-1)^k\frac{\zeta(k)}{k+1}z^{k+1}.
\end{equation}
\end{theorem}

\begin{proof}
Take the logarithm of the finite partial product in \cref{eq:Fdef} on
\(|z|<1\).  Each \(1+z/k\) lies in the disk centered at one of radius less than
one, so the Taylor branch is coherent.  Substitute \cref{eq:qseries}.  For
\(|z|\le\rho<1\),
\[
 \sum_{m=3}^{\infty}\sum_{k=1}^{\infty}
 \frac{|z|^m}{m k^{m-1}}
 \le 2\sum_{m=3}^{\infty}\frac{\rho^m}{m}<\infty.
\]
The finite rectangle identity is exact; the two tails have explicit geometric
bounds, so their difference is smaller than any requested dyadic radius.  The
linear and quadratic coefficients come from the prefactor.  The sign
\((-1)^{m-1}=(-1)^{m+1}\) is unchanged because the exponents differ by two.
\end{proof}

\subsection{Explicit truncation bounds}

Let \(T_M(z)\) contain the terms through degree \(M\), where \(M\ge3\).  If
\(|z|\le\rho<1\), then
\begin{equation}\label{eq:taylor-tail}
 \left|\log G(1+z)-T_M(z)\right|
 \le
 \frac{2\rho^{M+1}}{(M+1)(1-\rho)}.
\end{equation}
Indeed,
\[
 \sum_{m=M+1}^{\infty}
 \frac{\zeta(m-1)\rho^m}{m}
 \le \frac2{M+1}\sum_{m=M+1}^{\infty}\rho^m.
\]
For rational \(\rho=a/b<1\), the assertion that the right side is at most
\(2^{-p}\) is the natural-number inequality
\begin{equation}\label{eq:taylor-tail-cleared}
 2a^{M+1}b\,2^p
 \le (M+1)(b-a)b^{M+1}.
\end{equation}
The certificate producer increases \(M\) until this decidable inequality
holds.  Bounded search can be capped by a deliberately generous primitive-
recursive value, for example a polynomial-exponential function of
\(p+b\).  Soundness depends only on the checked inequality, not on optimality
of the cap.

The same argument gives derivative tails.  If \(0\le r\le M\), termwise
\(r\)-fold differentiation yields
\begin{equation}\label{eq:taylor-derivative-tail}
 \left|D^r\bigl(\log G(1+z)-T_M(z)\bigr)\right|
 \le
 \frac{2(M+1)^r\rho^{M+1-r}}
      {(1-\rho)^{r+1}}
\end{equation}
with a coarse but primitive-recursive constant.  A sharper falling-factorial
bound is available, but the coarse form is easier to compile and entirely
adequate.

\subsection{Exponentiating the series}

Exponentiating \cref{eq:taylor-wiki} gives
\begin{align}
 G(1+z)={}&(2\pi)^{z/2}
 \exp\!\left(-\frac{z+(1+\gamma)z^2}{2}\right)\notag\\
 &\times
 \exp\!\left(
 \sum_{k=2}^{\infty}(-1)^k\frac{\zeta(k)}{k+1}z^{k+1}
 \right).
 \label{eq:exp-taylor}
\end{align}
Comparing with the canonical product yields
\begin{equation}\label{eq:series-product-relation}
 \exp\!\left(
 \sum_{k=2}^{\infty}(-1)^k\frac{\zeta(k)}{k+1}z^{k+1}
 \right)
 =\prod_{k=1}^{\infty}E_k(z),
 \qquad |z|<1.
\end{equation}
This is not circular: both sides have independently supplied finite
approximants, and \cref{eq:q-double-sum} shows that their logarithmic
approximants differ by a rectangular tail tending effectively to zero.

For formalization, use the ball version
\[
 |u-v|\le2^{-p-c},\quad |u|,|v|\le B
 \Longrightarrow
 |\e^u-\e^v|\le \e^B2^{-p-c}.
\]
The exponential evaluator supplies an upper rational ball for \(\e^B\), and
the backward demand pass chooses \(c\).  No exact transcendental equality is
asked of the finite centers.

\subsection{The logarithmic derivative series}

Differentiate \cref{eq:taylor-logG}.  With
\[
 \phi(w):=\frac{d}{dw}\log G(w),
 \qquad
 \phi(1)=\frac{\log(2\pi)-1}{2},
\]
we obtain, for \(|w-1|<1\),
\begin{equation}\label{eq:phi-taylor}
 \phi(w)=\phi(1)-(1+\gamma)(w-1)
 +\sum_{k=2}^{\infty}(-1)^k\zeta(k)(w-1)^k.
\end{equation}
The tail on \(|w-1|\le\rho<1\) is bounded by
\begin{equation}\label{eq:phi-tail}
 \left|\sum_{k=M+1}^{\infty}(-1)^k\zeta(k)(w-1)^k\right|
 \le \frac{2\rho^{M+1}}{1-\rho}.
\end{equation}
These formulas prove the Taylor and logarithmic-derivative claims on the target
pages in a stronger complex-disk form.  One source states only \(0<z<1\);
that interval is a valid special case, not the full domain.

\subsection{Primitive-recursive names for the coefficients}

For each fixed integer \(s\ge2\), define
\[
 \zeta(s)=\sum_{n=1}^{N}\frac1{n^s}+R_{s,N},
 \qquad
 0<R_{s,N}\le\frac{1}{(s-1)N^{s-1}}.
\]
The finite sum is rational and primitive recursive.  Given \(p\), choose
\(N\) so that
\((s-1)N^{s-1}\ge2^{p+2}\).  Bounded powering and search produce a modulus.
Thus every coefficient in a fixed Taylor truncation is a fast rational name.
The development never requires a single object-language function that accepts
an arbitrary code for \(s\); the external theorem schema introduces the finite
list of fixed \(s\)'s needed at each truncation order.

Euler's constant has the standard name
\[
 \gamma_N=H_N-\log N,
 \qquad
 0<\gamma_N-\gamma<\frac1N,
\]
with a sharpened Euler--Maclaurin bound available from the Gamma packet.  The
constant \(\log(2\pi)\) is obtained from the named \(\pi\) and positive-log
evaluator.  Consequently every coefficient and every tail budget in
\cref{eq:taylor-logG} is expressible in the iPRA framework of
\cref{sec:ipra}.

\begin{roadmap}[Lean module \code{BarnesG/Taylor}]
Implement a finite coefficient function returning the rational balls for
\(\zeta(2),\ldots,\zeta(M-1)\), a Horner evaluator for \(T_M\), and the checker
for \cref{eq:taylor-tail-cleared}.  Prove the finite double-sum rearrangement by
bounded sum interchange, then use the canonical-factor tail to pass from the
finite rectangle to the product.  The object theorem is a ball inclusion, not
an equality of host-Lean real numbers.
\end{roadmap}
